% !TeX program = xelatex
% =============================================================
% 模板 03：中文通用课程作业（ctex + XeLaTeX + biblatex/biber）
% 编译： scripts\build.ps1 cn-homework
% 特点： 参考文献走 biber 后端（scripts\build.ps1 会自动识别 .bcf 并调用 biber）
% 版式： A4，上下 2.54cm，左右 3.17cm，正文小四，1.5 倍行距
% =============================================================
\documentclass[UTF8,a4paper,zihao=-4,fontset=windows]{ctexart}

\usepackage{geometry}
\geometry{top=2.54cm,bottom=2.54cm,left=3.17cm,right=3.17cm}
\usepackage{amsmath,amssymb,bm}
\usepackage{booktabs}
\usepackage{array}
\usepackage{graphicx}
\usepackage{caption}
\usepackage{enumitem}

% 参考文献：GB/T 7714-2015，biber 后端
\usepackage[backend=biber,style=gb7714-2015]{biblatex}
\addbibresource{refs.bib}

\usepackage[hidelinks]{hyperref}

\linespread{1.5}
\pagestyle{plain}   % 页码居中页脚，不显示页眉
\captionsetup{font=small,labelsep=quad}
\captionsetup[table]{position=top}

% 需要代码块 / 算法伪代码 / 子图 / 跨页长表 / 图表按节编号时：
% 见 references\latex-blocks.md，复制对应导言区与正文块即可

\ctexset{
  section    = {format = \zihao{4}\heiti},
  subsection = {format = \zihao{-4}\heiti}
}

\begin{document}

% ------------------------- 标题区（改这里） -------------------------
\begin{center}
  {\zihao{3}\heiti 作业标题占位：第二章 电路分析方法 习题解答}\\[0.6em]
  {\zihao{5} 张三\quad 2026000000\quad 电气工程及其自动化 2 班}
\end{center}

\vspace{0.4em}\hrule\vspace{0.8em}

\section{题目}

占位题目：如图（此处省略插图）所示电路，已知 $R_1=10\,\Omega$、$R_2=20\,\Omega$、$U_S=12\,\text{V}$，求各支路电流。

\section{解题思路}

采用节点电压法。以电路下端为参考节点，则上端节点电压 $U_n$ 满足
\begin{equation}
  \frac{U_n-U_S}{R_1}+\frac{U_n}{R_2}=0,
  \label{eq:kcl}
\end{equation}
整理得
\begin{equation}
  U_n=\frac{U_S R_2}{R_1+R_2}.
  \label{eq:solve}
\end{equation}

\section{计算过程}

将 $U_S=12\,\text{V}$、$R_1=10\,\Omega$、$R_2=20\,\Omega$ 代入式~\eqref{eq:solve}，得 $U_n=8\,\text{V}$，各支路电流见表~\ref{tab:result}。

\begin{table}[htbp]
  \centering
  \caption{各支路电流计算结果}
  \label{tab:result}
  \begin{tabular}{lccc}
    \toprule
    支路 & 电阻/$\Omega$ & 电压/V & 电流/A \\
    \midrule
    $R_1$ 支路 & 10 & 4  & 0.400 \\
    $R_2$ 支路 & 20 & 8  & 0.400 \\
    \bottomrule
  \end{tabular}
\end{table}

\section{结果讨论}

计算结果表明两支路电流相等，与串联分压的结论一致\cite{zhang2024demo}；若考虑电源内阻，可用戴维南定理修正，方法见文献\cite{smith2023demo}。

\section{小结}

本题用节点电压法求解，核心是把基尔霍夫电流定律写成式~\eqref{eq:kcl} 的形式，再代数求解即可。

\printbibliography[title={参考文献}]

\end{document}
